\section{Quadratic wake tests for non-axis drift}\label{app:quadratic}
The rule is white/right or black/left, followed by toggling the current site
and moving one edge. The black field is a complete finite subset of $\Z^2$.
Suppose a nonempty turn word of length $P$ repeats forever, restores the heading,
and has nonzero full displacement $d$. Finite preludes and genuinely unread
debris are allowed. We do not replace $d$ by its primitive direction.

The stable-order and port-factorization method of~\cite{Jillhewar2026}
lets us eliminate the unknown event graph and test the wake directly.
Assume $d=(a,b)$, $ab\ne0$. A half-turn makes $a>0$; heading restoration implies
$a+b$ even. Let $O$ consist of full $d$-classes with odd phase multiplicity and
write $g=|O|$. Canonical representatives are
\[
 \pi_d(x,y)=(r,y-b\lfloor x/a\rfloor),\qquad r=x\bmod a,\quad 0\le r<a.
\]
Checker parity $c(r,y)=(r+y)\bmod2$ is well-defined on the quotient. Initially
use $Q=\{c=0\}$ for vertical-arrival vertices and $C=\{c=1\}$ for horizontal
arrivals. Lemma~\ref{qw:lem:cancel} removes dependence on this convention.

\subsection{The graph reconstructed from a wake}
Suppose that every horizontal and vertical quotient component contains an even
number of $O$ vertices. On each vertical line $r=x\bmod a$, sort the canonical
$y$ coordinates and pair consecutive endpoints by finite intervals. Horizontal
components have representative row $Y=y\bmod |b|$. Express a point there using
\[
 k=(y-Y)/b,\qquad X=x-ka,
\]
then sort and pair the $X$ coordinates. Let $G_d(O)$ be the union of these
horizontal and vertical interval edges. Orient horizontal edges from $Q$ to
$C$ and vertical edges from $C$ to $Q$.

Order vertices by $(y,r)$ and ports negative before positive; positive means
movement $E$ or $N$. An edge has source and target keys $(\text{vertex},\text{port})$.
Let $E_d(O)$ count, modulo two, edge pairs whose source and target orders disagree.
Let $N_+(O)$ count positive retained edges, and put $C(O)=|O\cap C|$. Define
\begin{align}
 \qd(O)&=E_d(O)+N_+(O)+C(O)\pmod2,\label{qw:eq:character}\\
 \Od(O)&=\sum_{(r,y)\in O}\left\lfloor y/|b|\right\rfloor\pmod2.
 \label{qw:eq:winding}
\end{align}
The finite graph may be enormous, but Section~\ref{qw:sec:algorithm} does not
construct its individual edges.

\begin{theorem}[Wake obstruction]\label{qw:thm:main}
Every finite-support periodic realization with non-axis full drift satisfies
\[
 \boxed{\Od(O)=1,\qquad \qd(O)=1+a\pmod2.}
\]
There is no restriction on period, transverse width, growth, retreat, seed size,
or wake shape. The characters are independent of representatives, translation,
phase origin and checker/heading-axis convention; no heading is a screen input.
\end{theorem}

\begin{theorem}[Quadratic structure]\label{qw:thm:quadratic}
On the $\F$ vector space of finite indicators with even sums in every horizontal
and vertical quotient component, $\qd$ is a quadratic form. In particular
$\qd(A+B)+\qd(A)+\qd(B)$ is bilinear. Addition is symmetric difference of
indicators, not superposition of physical trajectories.
\end{theorem}

\begin{theorem}[Compressed complexity]\label{qw:thm:complexity}
Given $d$ and explicit representatives of $O$ in binary, signed-residue
prefilters followed by Theorem~\ref{qw:thm:main} are computable in polynomial bit
time. After the density filter, at most
$R=g(a+1)\le g(g/2+1)$ arithmetic edge runs suffice, and the evaluator uses
$O(g^4\log(g+1))$ integer arithmetic operations on $O(B+\log(g+1))$-bit integers,
where $B$ bounds coordinate and drift bit lengths. Passing means only
\texttt{NOT\_EXCLUDED}, not realizability.
\end{theorem}

\subsection{Boundary reconstruction and winding}
In a full $d$-class write phase positions $q_i=b_v+k_i d$. A sufficiently distant
cell $b_v+nd$ is initially white and eventually has all the phase visits, at
times $nP+i-Pk_i$. The necessary stable-class condition of~\cite{Jillhewar2026} therefore
orders visits by $\tau_i=i-Pk_i$ and gives letters $R,L,R,L,\ldots$.
The signed excess in a class is its odd indicator $D_v\in\{0,1\}$.
Thus $g=\sum D_v$, and heading restoration gives $4\mid g$.

The checker/heading-axis invariant assigns incoming axes alternately to the
two checker colours. Every visit uses one horizontal and one vertical edge.
Modulo two, the horizontal incident-edge boundary and the vertical
incident-edge boundary are consequently both $D$.

Because $b\ne0$, the horizontal quotient graph is a union of infinite paths:
no nonzero horizontal vector is in $\Z d$. The analogous fact holds vertically
because $a\ne0$. On an infinite path, a finite edge set with prescribed boundary
is unique: its indicator is the cumulative endpoint parity, hence exactly the
intervals pairing consecutive endpoints. A nonempty finite boundary-free edge
set on a line cannot exist. Therefore the actual edge-count parities are
precisely $G_d(O)$, whatever their unknown even parts may be.

This argument explains the domain restriction. For axis drift, one direction
has quotient cycles, and a boundary-free cycle can be added to parity data.
See Section~\ref{sec:limitations} for the resulting axis-scope boundary.

The actual period has signed seam crossing number one across the cut from
$x$ residue $a-1$ to $0$. For a horizontal interval from $X_L$ to $X_R$, its
seam crossing parity is $\lfloor X_R/a\rfloor-\lfloor X_L/a\rfloor$.
Summing over paired intervals equals the sum of these floors over all endpoints
modulo two. A canonical endpoint $(r,y)$ has $X=r-ak$, with
$k=(y-(y\bmod|b|))/b$, so $\lfloor X/a\rfloor=-k$ has parity
$\lfloor y/|b|\rfloor$. This proves~\eqref{qw:eq:winding} and $\Od=1$.
On arbitrary representatives the same weight is
$\lfloor x/a\rfloor+\lfloor y/|b|\rfloor\pmod2$; adding $d$ changes it by the
even integer $1+\operatorname{sign}(b)$.

\subsection{The one-tour sign and its wake reduction}
At a vertex order its $m$ events by stable rank $k=0,\ldots,m-1$.
Let $\alpha_k,\beta_k$ be positive incoming and outgoing port bits, with totals
$A,B$, and put $p_v=A+B\pmod2$. At $Q$, a left turn means unequal bits;
at $C$, it means equal bits. For any binary list,
\begin{equation}
 \inv(\alpha)=A(m-1)+\sum_k k\alpha_k+\binom A2\pmod2.
 \label{qw:eq:binary}
\end{equation}
Indeed, each one at position $k$ sees $m-1-k$ later positions, from which its
later ones are subtracted. Summing the latter gives $\binom A2$.

Put $K=\sum_{L\text{ events}}\text{stable rank}\pmod2$.
Equation~\eqref{qw:eq:binary} applied to both port lists leaves local corrections
$p_v(m-1)$ at $Q$ and $p_v(m-1)+\binom m2$ at $C$, besides port inversions and
positive-edge binomials. At $Q$, $m\equiv D_v\pmod2$. At $C$, counting unequal
bits gives
\[
 D_v=2(A+B)-4\sum_k\alpha_k\beta_k-m,
 \qquad m\equiv2p_v-D_v\pmod4.
\]
For $D_v=0,1$ the two corrections are respectively $p_v(D_v+1)$ and
$D_v(p_v+1)$. The positive-edge binomials cancel between source and target.
If $I$ is the sum of all incoming/outgoing port-list inversion parities,
\begin{equation}
 K=I+L,\qquad L=\sum_Qp_v(D_v+1)+\sum_CD_v(p_v+1).\label{qw:eq:local}
\end{equation}

Along a fixed directed quotient edge of voltage $\nu$, target stable time is
source stable time plus $1-P\nu$. The same constant applies to every occurrence,
including the cyclic seam. Occurrences therefore match first-to-first.
Sorting outgoing ports, following edge blocks and unsorting incoming ports
factors the phase successor. This successor is \emph{one} even-length $P$-cycle,
of odd sign; hence
\begin{equation}
 I+E=1,\qquad K=1+E+L.\label{qw:eq:factor}
\end{equation}
A block reversal of sizes $u,v$ contributes $uv\pmod2$. Thus $E$ depends only
on edge-count parities and equals $E_d(O)$. This is the general version of the
attributed special-drift factorization, not a finite-period extrapolation.

\begin{lemma}[Cancellation and conventions]\label{qw:lem:cancel}
For $G_d(O)$ one has
\[
 \sum_{v\in O}p_v=0,\qquad L=N_+(O)+C(O)\pmod2.
\]
The character is unchanged by the total vertex order or by exchanging the
checker/heading-axis convention.
\end{lemma}
\begin{proof}
The positive incidence of a paired horizontal interval at its two endpoints
has parity $1+c(\text{left})+c(\text{right})$. A vertical interval has the same
formula. There are $g/2$ intervals of each kind; each endpoint checker bit is
counted twice. The total is $g=0\pmod2$.
Also $\sum_Qp_v=N_+$, since every positive edge touches exactly one $Q$ vertex.
Expanding~\eqref{qw:eq:local} proves the stated expression for $L$.

Swapping vertices $v,w$ in the order changes $E$ by
$\operatorname{out}(v)\operatorname{out}(w)+
\operatorname{in}(v)\operatorname{in}(w)=2D_vD_w=0\pmod2$.
Exchanging checker conventions reverses all edges and complements the ports.
Inversion sign is invariant under permutation inversion; the extra port-block
reversals contribute $\sum_v(1+D_v)p_v=0$, by endpoint cancellation and
$\sum_vp_v=0$. The edge set has even size, because the sum of outgoing degree
parities is $g$. Thus replacing $N_+$ by $|G|-N_+$ does not change its parity;
neither does replacing $C$ by $g-C$.
\end{proof}
Translations merely relabel vertices and possibly exchange conventions, so
$\qd$ is translation and phase invariant. Moving a seam does not change winding,
because every fibre population is even.

For an actual stable class word, every $L$ has odd rank. Hence
$K=\#L=(P-g)/2\pmod2$. There are $P/2$ horizontal moves with net displacement
$a$, so $P/2\equiv a\pmod2$; since $4\mid g$, $K\equiv a$.
Equations~\eqref{qw:eq:character}, \eqref{qw:eq:local}, \eqref{qw:eq:factor} prove
Theorem~\ref{qw:thm:main}.

The reconstruction map $D\mapsto G_d(D)$ is linear by uniqueness of a finite
boundary antiderivative. The inversion parity is quadratic in edge indicators,
and $N_+,C$ are linear. This proves Theorem~\ref{qw:thm:quadratic}.
For $(2,-2)$, the previous evaluation in Appendix~\ref{app:mixed} is
$\qd=J+V/2$ under signed residues, where $J$ counts $Q$-before-$C$ odd-class pairs
in transverse order and $V=|O_Q|$. Thus the earlier formula $J+V/2=1$ is recovered.

\subsection{Compressed evaluation}\label{qw:sec:algorithm}
The signed residue theorem~\cite[Section~7, ``Signed mod-four wake-residue identity'']{Jillhewar2026} gives, in every $x$ residue modulo $a$ and
every $y$ residue modulo $|b|$, signed odd-class sum $2\pmod4$, with sign
$(-1)^{x+y}$. In particular $g\ge2\max(a,|b|)$.
Check positivity and $4\mid g$, then density \emph{before} allocating residue-sized
arrays; rejected inputs need no graph construction. Both drift magnitudes are
then at most $g/2$. Signed sums also enforce even fibre populations.

An arithmetic edge run has $n$ edges with source and target keys
\begin{equation}
 (s_y+hi,s_x,p),\quad (t_y+hi,t_x,p),\qquad 0\le i<n,\quad h>0.
 \label{qw:eq:run}
\end{equation}
A vertical interval splits by bottom-site parity into at most two step-two runs.
For a horizontal interval on representative row $Y$, split its left positions
$X=r+ak$ by $r$, and, for odd $a$, by $k\bmod2$. Canonical heights change by
$-b$ per $k$. Reverse the index if necessary to obtain positive step $|b|$ or
$2|b|$ and constant port bits. There are at most $2a$ runs per horizontal interval,
so $R\le g(a+1)$. All runs have disjoint edges.

A run has no internal inversions. Between runs $A,B$, set
\[
 z=s_A+h_Ai-s_B-h_Bj,\qquad
 \delta=(t_A-s_A)-(t_B-s_B).
\]
Let $s=0$ if the source minor key $(x_A,p_A)$ precedes $(x_B,p_B)$, otherwise
$s=-1$. Define $u=-\delta$ plus $0$ or $-1$ by the analogous target comparison.
The edge pair is inverted exactly when
\begin{equation}
 \min(s,u)<z\le\max(s,u).\label{qw:eq:inversion-range}
\end{equation}
Two cumulative pair counts suffice. For
$ui-vj\le T$, $0\le i<n$, $0\le j<m$, put
\[
 \ell=\clip(\lceil-T/v\rceil,0,m),\qquad
 h=\clip(\lceil(un-T)/v\rceil,0,m).
\]
The exact number of pairs is
\begin{equation}
 (m-h)n+(h-\ell)+
 \sum_{j=0}^{h-\ell-1}\left\lfloor\frac{vj+T+v\ell}{u}\right\rfloor.
 \label{qw:eq:pairs}
\end{equation}
There are no eligible $i$ below $\ell$, all $n$ are eligible above $h$, and the
middle contributes $\lfloor(T+vj)/u\rfloor+1$.

The floor sum $S(n,m,a,b)=\sum_{i=0}^{n-1}\lfloor(ai+b)/m\rfloor$ uses Euclidean
lattice transposition. Remove quotient parts of $a,b$, contributing
$\lfloor a/m\rfloor n(n-1)/2+\lfloor b/m\rfloor n$.
For $0\le a,b<m$ the remainder is zero if $an+b<m$; otherwise, with
$n'=\lfloor(an+b)/m\rfloor$ and $b'=(an+b)\bmod m$,
\begin{equation}
 S(n,m,a,b)=S(n',a,m,b').\label{qw:eq:euclid}
\end{equation}
At positive height $j$ there are $n-\lceil(mj-b)/a\rceil$ eligible indices.
Reindex heights from $n'$ downwards, using $mn'-b=an-b'$, to obtain the right
side of~\eqref{qw:eq:euclid}. Division first handles negative offsets exactly.
Euclidean reduction needs logarithmically many stages.

All run-pair counts have polynomial bit length even for runs of $10^{100}$ edges.
Summing positive run lengths gives $N_+$; $C$ and $\Od$ use only the endpoints.
The run bound, $R^2$ run pairs, and~\eqref{qw:eq:pairs}--\eqref{qw:eq:euclid} prove
Theorem~\ref{qw:thm:complexity}. A passing wake proceeds to the chronology-exact inverse synthesis of
Appendix~\ref{app:general}.

\subsection{An unbounded growth-eight exclusion}\label{qw:sec:family}
Let $d=(3,1)$ and
\begin{align*}
 O_* =\{&(0,1),(0,3),(1,-2),(1,-1),\\[-2pt]
        &(1,0),(1,2),(2,0),(2,2)\}.
\end{align*}
For $n\ge1$, put $\lambda=12n+1$ and let $O_n$ be the classes of $\lambda q$,
$q\in O_*$. They are eight distinct classes with canonical representatives
\[
 (r,(12y-4r)n+y),\qquad (r,y)\in O_*.
\]
Their signed $x$-fibre sums are $(-2,-2,2)$ and their sole $y$-fibre sum is $-2$.
They have growth eight, pass density, and have winding
$\sum\text{canonical }y=28n+5\equiv1\pmod2$.

\begin{corollary}\label{qw:cor:family}
No finite-support periodic highway of drift $(3,1)$ has wake $O_n$, for any
integer $n\ge1$ and any period. Their algebraic transverse span, measured by
$T(x,y)=3y-x$, is $16(12n+1)$ and is unbounded.
\end{corollary}
\begin{proof}
The affine runs in Appendix~\ref{qw:app:runs} give
\[
 |G_d(O_n)|=192n+16,\quad N_+(O_n)=96n+9,\quad
 \#\text{inversions}=400n+19,\quad C(O_n)=5.
\]
Thus $\qd(O_n)=496n+33\equiv1$, whereas Theorem~\ref{qw:thm:main} requires
$1+a\equiv0$. For completeness, the calculation is an all-parameter integer
identity, not an interpolation. If $L<R$ are horizontal representative-row
endpoints before dilation, the index bounds at residue $r$ are
\[
 k_{\min}=4Ln+\lceil(L-r)/3\rceil,\qquad
 k_{\max}=4Rn+\lfloor(R-1-r)/3\rfloor.
\]
Splitting parity gives 32 runs of step two and positive affine lengths on
$n\ge1$. For a parity-compatible inversion difference $z$, set
$k=(z-s_A+s_B)/2$. Its count is exactly
\begin{equation}
 \max\{0,\min(n_B,n_A-k)-\max(0,-k)\}.\label{qw:eq:affine-pair}
\end{equation}
Every minimum/maximum branch is verified uniformly: an affine expression
$An+B$ is nonnegative for every $n\ge1$ exactly when $A\ge0$ and $A+B\ge0$.
Summing the 288 parity-compatible terms among 496 run pairs gives $400n+19$.
The supplement implements this symbolic arithmetic and fails rather than
extrapolating if a comparison changes sign on the parameter domain.
\end{proof}

The unscaled $O_*$ also fails: its counts are 16 retained edges, 9 positive
edges, 23 inversions and $C=5$. For contrast,
\[
 \{(0,1),(0,5),(1,-4),(1,-2),(1,5),(1,6),(2,-4),(2,4)\}
\]
has the same signed fibre populations and winding, but $\qd=0$.
It satisfies this necessary parity condition. Thus the character strengthens
the specified signed-residue, density and winding conditions. Physical
realization additionally requires chronological compatibility of the visits.

\subsection{Exact computational checks}
The proof and complexity theorem do not depend on bounded experiments. Exact
affine algebra establishes the entire family parameter range. Independent implementations
check 10,000 floor sums and 10,000 pair-count instances against literal sums,
4,000 individual-edge reconstructions for eight drift types, vertex-order and
checker-convention independence, and quadratic polarization.
All 8,192 closed masks in an eight-layer/two-lane window reproduce in Appendix~\ref{app:mixed}.

An independent C++ vector-turn implementation enumerates all 1,398,100 binary
words of even length at most 20 and checks the unspecialized port factorization
on all 456,120 non-axis heading-restoring words. The 84 words with class balances
in $\{0,1\}$ are additionally cross-checked against the wake-only formula. These
geometric words are not asserted physically realizable.
All 104 standard-highway phases, four rotations and powers one through four
provide 1,664 positive controls using their \emph{full} displacements.
Representatives with large coordinates and $n=10^{100}$
are checked without traversing their gaps; the latter still uses 32 runs.
Exact reports and source manifests are in the supplement.

\subsection{Affine edge-run certificate}\label{qw:app:runs}
All runs have step two and index $i=0,\ldots,\ell-1$. A row gives length $\ell$,
source $(s_x,s_y+2i)$, target $(t_x,s_y+\delta+2i)$ and port $p$.
These are edges of the parity graph, not time intervals of a proposed ant orbit.
The row order is only an enumeration of runs; inversions use vertex order $(y,r)$.
Formula~\eqref{qw:eq:affine-pair} on these rows yields the counts in
Corollary~\ref{qw:cor:family} for the whole domain $n\ge1$.

\begin{longtable}{r r r r r r r}
\toprule
Run & $\ell$ & $s_x$ & $s_y$ & $t_x$ & $\delta$ & $p$\\
\midrule
\endhead
1&$8n+1$&0&$20n+2$&1&0&1\\
2&$8n+1$&1&$20n+3$&0&0&0\\
3&$8n$&2&$20n+4$&1&0&0\\
4&$8n+1$&1&$20n+3$&2&0&1\\
5&$8n$&2&$20n+4$&0&-1&1\\
6&$8n+1$&0&$20n+2$&2&1&0\\
7&$2n$&0&$12n+2$&1&0&1\\
8&$2n$&1&$12n+3$&0&0&0\\
9&$2n$&2&$12n+2$&1&0&0\\
10&$2n$&1&$12n+3$&2&0&1\\
11&$2n+1$&2&$12n+2$&0&-1&1\\
12&$2n$&0&$12n+2$&2&1&0\\
13&$2n$&0&$-8n$&1&0&1\\
14&$2n$&1&$-8n+1$&0&0&0\\
15&$2n+1$&2&$-8n$&1&0&0\\
16&$2n$&1&$-8n+1$&2&0&1\\
17&$2n$&2&$-8n+2$&0&-1&1\\
18&$2n$&0&$-8n$&2&1&0\\
19&$6n+1$&0&$-28n-2$&1&0&1\\
20&$6n$&1&$-28n-1$&0&0&0\\
21&$6n$&2&$-28n$&1&0&0\\
22&$6n+1$&1&$-28n-1$&2&0&1\\
23&$6n$&2&$-28n$&0&-1&1\\
24&$6n+1$&0&$-28n-2$&2&1&0\\
25&$12n+1$&0&$12n+3$&0&-1&0\\
26&$12n+1$&0&$12n+1$&0&1&1\\
27&$6n+1$&1&$-28n-2$&1&1&1\\
28&$6n$&1&$-28n$&1&-1&0\\
29&$12n+1$&1&$-4n$&1&1&1\\
30&$12n+1$&1&$-4n+2$&1&-1&0\\
31&$12n+1$&2&$-8n+1$&2&-1&0\\
32&$12n+1$&2&$-8n+1$&2&1&1\\
\bottomrule
\end{longtable}
